% Candidate S44. Author /root; GPT-6 Astra Pro; effort not exposed.
% Date 2026-11-21. Author and independent semireg entire-proof readings PASS.
\begin{theorem}[S-44: Embedded compact ball surfaces in abelian fourfolds are not Bloch semiregular]
\label{thm:ball-fourfold-normal-obstruction}
Let $Y$ be a smooth compact quotient of the complex two-ball by a
fixed-point-free discrete group, and let $i:Y\hookrightarrow A$ be a
closed immersion into a complex abelian fourfold. Write
$\chi=\chi(Y,\mathcal O_Y)$ and $N=N_{Y/A}$. Then the image generates
$A$, and
\[
 \chi\ge3,\qquad h^0(Y,N)=h^2(Y,N)=4,
 \qquad h^1(Y,N)=8+4\chi.
 \tag{S44.1}
\]
The classical Bloch semiregularity map therefore satisfies
\[
 \sigma_i:H^1(Y,N)\longrightarrow H^3(A,\Omega_A^1),
 \qquad \dim\ker\sigma_i\ge4\chi-8\ge4.
 \tag{S44.2}
\]
In particular no such embedding is Bloch semiregular. The assertion
requires neither a CM action nor a specified four-factor period.

If a compact ball-quotient surface maps finitely and etale onto a smooth
reduced surface $Z\subset A$, the same conclusion applies to the
embedding of $Z$. It does not apply merely from a generically finite
ramified map onto $Z$, and does not assert that any 032 image is smooth
or that its map is unramified.
\end{theorem}
\begin{proof}
\noindent\textbf{The ball inputs and generation.}
We use Hirzebruch proportionality and Calabi--Vesentini vanishing in
their stated compact, torsion-free setting. Their exact primary
locators are listed after the proof. They give
\[
 K_Y^2=9\chi,\qquad c_2(Y)=3\chi,\qquad \chi>0,
 \qquad H^0(Y,T_Y)=H^1(Y,T_Y)=0.
 \tag{S44.3}
\]
The canonical bundle is ample. We also use directly the elementary
hyperbolicity consequence of the given universal cover: there is no
nonconstant map to $Y$ from a curve of genus zero or one. Indeed, after
passing to the universal cover of the curve such a map lifts to the
bounded domain $\mathbb B^2$; the lift is constant by the maximum
principle in genus zero and by Liouville's theorem in genus one.
Thus the normalization of every curve in $Y$ has genus at least two.

Let $A_0$ be the smallest abelian subvariety whose translate contains
$i(Y)$. Its dimension is at least two. If it is two, the image is that
entire translate, which is not a ball quotient. If its dimension is
three, $Y$ is a smooth divisor in a smooth abelian threefold.
Adjunction and the tangent sequence give
$c(T_Y)=(1+c_1(\mathcal O(Y)|_Y))^{-1}$, and hence
$K_Y^2=c_2(Y)$, contradicting (S44.3). Therefore $A_0=A$.
The induced Albanese homomorphism is surjective, so pullback of
holomorphic one-forms is injective. In particular
\[
 q:=h^1(Y,\mathcal O_Y)=h^0(Y,\Omega_Y^1)\ge4.
 \tag{S44.4}
\]

\noindent\textbf{The exact normal cohomology.}
The normal sequence and (S44.3) give
\[
 0\longrightarrow T_Y\longrightarrow\mathcal O_Y^{\oplus4}
   \longrightarrow N\longrightarrow0,
 \qquad h^0(N)=4,
 \qquad H^1(\mathcal O_Y)^{\oplus4}\hookrightarrow H^1(N).
 \tag{S44.5}
\]
Since $N$ has rank two, its determinant is $K_Y$, and
\[
 c_1(N)=K_Y,\qquad
 c_2(N)=K_Y^2-c_2(Y)=6\chi.
 \tag{S44.6}
\]
The rank-two identity $N^*\otimes\det N\simeq N$ and Serre duality
imply $h^2(N)=h^0(N)=4$. Riemann--Roch for this rank-two bundle gives
\[
 \begin{aligned}
 \chi(Y,N)
 &=2\chi+\tfrac12 c_1(N)\bigl(c_1(N)-K_Y\bigr)-c_2(N)\\
 &=-4\chi.
 \end{aligned}
 \tag{S44.7}
\]
It follows that $h^1(N)=8+4\chi$. The injection in (S44.5) also gives
\[
 4q\le8+4\chi,\qquad q\le\chi+2.
 \tag{S44.8}
\]
Together with (S44.4) and positivity of the integral $\chi$, this
already proves $\chi\ge2$.

\noindent\textbf{Eliminating the boundary $\chi=2$.}
Suppose now that $\chi=2$. Equations (S44.4) and (S44.8) give $q=4$,
and then $p_g=\chi+q-1=5$. All holomorphic one-forms on $Y$ are
therefore pulled back from $A$. The map
\[
 \bigwedge^2H^0(A,\Omega_A^1)\longrightarrow H^0(Y,K_Y)
 \tag{S44.9}
\]
has a six-dimensional source and a five-dimensional target, so has
a nonzero kernel. Choose a nonzero invariant two-form $\omega$ in
that kernel. Its rank on the four-dimensional tangent space of $A$
is either two or four.

If the rank is four, $\omega$ is holomorphic symplectic and $Y$ is a
Lagrangian submanifold: its two-dimensional tangent spaces are
isotropic because $i^*\omega=0$. Pairing with $\omega$ gives an
isomorphism of vector bundles
\[
 N\simeq\Omega_Y^1.
\]
Then $c_2(N)=c_2(Y)$, whereas (S44.6) gives
$c_2(N)=K_Y^2-c_2(Y)$. This forces $K_Y^2=2c_2(Y)$, in contradiction
with (S44.3) and $c_2(Y)>0$.

Thus $\omega$ has rank two and is decomposable, say
$\omega=\alpha\wedge\beta$ for linearly independent one-forms on
$A$. Their restrictions remain linearly independent by generation.
The Castelnuovo--de Franchis theorem now gives a surjective
holomorphic map with connected fibres
\[
 f:Y\longrightarrow C,\qquad b=g(C)\ge2.
 \tag{S44.10}
\]
Let $g$ be the genus of a smooth fibre. The hyperbolicity observation
above gives $g\ge2$. The Zeuthen--Segre formula for a fibration of a
smooth compact surface gives
\[
 e(Y)=4(b-1)(g-1)+\mu,\qquad \mu\ge0.
 \tag{S44.11}
\]
Here $e(Y)=c_2(Y)=6$. The positive integers $b-1$ and $g-1$
therefore have product at most $3/2$, so $b=g=2$.

Every fibre of this genus-two fibration is in fact smooth and reduced.
To see this without assuming semistability, write a fibre divisor as
$F=\sum_j m_jC_j$, with distinct integral curves $C_j$ and
$m_j\ge1$. Since $F\cdot C_j=0$ and distinct curves have
nonnegative intersection, $C_j^2\le0$. The normalization of each
$C_j$ has genus at least two, so its arithmetic genus is at least
two. Adjunction therefore gives
\[
 K_Y\cdot C_j=2p_a(C_j)-2-C_j^2\ge2.
 \tag{S44.12}
\]
A smooth fibre has genus two and square zero, so numerical equivalence
of fibres gives $K_Y\cdot F=2$. Consequently
\[
 2=\sum_jm_j K_Y\cdot C_j\ge2\sum_jm_j.
\]
There is exactly one component with multiplicity one. For that curve
$C$, one has $C^2=0$ and $p_a(C)=2$. Its normalization has genus at
least two, hence its total delta invariant is zero and $C$ is smooth.
Thus every fibre is smooth. The smooth fibration is differentiably
locally trivial, and multiplicativity of the topological Euler
characteristic gives
\[
 e(Y)=e(C)e(F)=(2-4)^2=4,
\]
contradicting $e(Y)=6$. This excludes $\chi=2$ and proves (S44.1).

\noindent\textbf{The Bloch obstruction and the etale-image scope.}
For a codimension-two embedding the classical semiregularity target
is $H^3(A,\Omega_A^1)$. The cotangent bundle of an abelian fourfold
is trivial, so this target has dimension
$4\binom43=16$. Equation (S44.1) and rank-nullity give (S44.2).

Finally, a connected finite etale quotient of a compact ball quotient
has the same simply connected universal covering domain. A smooth
surface $Z$ as in the statement is therefore itself a compact ball
quotient, so the theorem applies to $Z\hookrightarrow A$. A ramified
map does not have this consequence, and no such extension is used.
\end{proof}

\paragraph{Exact imported sources and actual reading scope.}
Hirzebruch, \emph{Automorphe Formen und der Satz von Riemann--Roch}
(1958), Section~3, Satz~4, printed p.~137, with the compact dual of
the ball identified as projective space in Section~5, printed p.~143.
The statement and its preceding proportionality argument were read
in the author's MPIM collection; the positive canonical-bundle
statement is in Section~3, printed pp.~135--136. Source:
\url{https://hirzebruch.mpim-bonn.mpg.de/id/eprint/89/}.
This is the imported proportionality theorem, not a new proof of it.

Calabi--Vesentini, \emph{Sur les varietes complexes compactes localement
symetriques}, Bull. Soc. Math. France~87 (1959), Theoreme~3,
printed p.~315, gives both $H^0(T_Y)=0$ and $H^1(T_Y)=0$ for the
ball of dimension two (type $I_{1,2}$). The whole theorem and its
hypotheses were read; the announcement explicitly sends its full
proof elsewhere. Source:
\url{https://www.numdam.org/article/BSMF_1959__87__311_0.pdf}.

The classical Castelnuovo--de Franchis statement is given explicitly
by Catanese, \emph{Fibred Kahler and quasi-projective groups},
Advances in Geometry, Special Issue (2003), Section~1, printed p.~S13;
an exact numbered version is Pierre Py, \emph{Lectures on Kahler
groups}, Theorem~4.1, printed p.~46 in the author manuscript. Its full
statement, including connected fibres, was retrieved and read; no
independent audit of its full proof is claimed. Catanese's Section~6,
proof of Theorem~6.3, printed p.~S22, explicitly states the
Zeuthen--Segre formula (S44.11) and nonnegativity for a smooth compact
surface fibration. That entire short proof was read. Sources:
\url{https://www.mathe8.uni-bayreuth.de/de/team/prof-fabrizio-catanese-old/pdf/088.pdf}
and
\url{https://www-fourier.ujf-grenoble.fr/~py/Documents/Livre-Kahler/lectures.pdf}.

The normal-sequence, Serre-duality and surface Riemann--Roch inputs
are the standard intersection-theoretic statements used elsewhere in
this paper (Fulton, Theorem~18.3 and Example~18.3.11 for the
Riemann--Roch conventions). The Bloch target is the classical one
specified and sourced in S-02. The new argument is the combination
of these inputs, the exact normal dimensions, and the elimination
of the small boundary case. Its framework placement is FGA,
Bourbaki expose~221, Lemme~2.4, p.~255 and Theoreme~3.2, p.~260,
for embedded deformation spaces, and SGA~6, Expose~0, Section~1,
equations (1.1)--(1.3), p.~1, for the Chern-character intersection
framework. No ball-quotient or semiregularity theorem is attributed
to those Grothendieck locators.
